\cotefile{34}
\scanpage{149}
\mspage{V.21}
% archive sheet number: 148

5. \uline{Application des th. de Lefschetz local} (SGA 2 $\overline{\underline{\unsure{\mathrm{XIV}|\mathrm{XV}}}}$ 4).

5.0. Cf. les commentaires dans \unsure{3.3|3.2}. pour les hypothèses faites quand nous \unsure{utiliserons} les théorèmes de Lefschetz locaux.

5.1. Soit $f:X\longrightarrow S$ loc. de type fini, $S$ str. loc. \unsure{hensélien}, et soit $x\in X_0$. Soit $m'\in\mathbb{Z}$ et supposons que l'on \unsure{sache} \deleted{(à titre d'hypothèse}
\[
({\unsure{\underline{\Phi}|\underline{\varphi}}}^i_{\unsure{\nu|v}})_{\bar{x}'}=0
\quad\text{pour tt \unsure{gén. stricte} de $x$ dans $X_0$, et $i\leq m'$.}
\tag*{\tikz[baseline=(n.base)]{\node[inner sep=1pt] (n) {5.1.1};\draw (n.center) ellipse (0.62cm and 0.22cm);}}
\]

\deleted{On} Nous \unsure{considérons} cette hypothèse comme une hypothèse de \unsure{récurrence|voisinage}, et nous proposons de donner des conditions moyennant \deleted{\unsure{un}} \unsure{quoi} on conclura
\[
(\underline{\Phi}^{i})_{\bar{x}}=0\quad\text{si}\quad i\leq m.
\tag*{(*)}
\]

Notons que l'on sait déjà que (5.1.1) implique
\[
\deleted{H^i(\bar{V},\Lambda_\nu)\xrightarrow{\sim}H^i(\unsure{\widetilde{X}_{\bar{\eta}}|\widetilde{X}_{\bar{y}}},\Lambda_\nu)\quad\text{si}\quad i<m'}
\tag*{\deleted{(5.1.1)}}
\]
\deleted{\unsure{en vertu de}}
\[
(\underline{\Phi}^i_\nu)_{\bar{x}}=H^i_x(\widetilde{X},\Lambda_\nu)
=\varinjlim_{\unsure{h|k}} H^i_x(X(h),\Lambda_\nu)
\quad\text{pour}\quad i\leq m'.
\tag*{(5.1.2)}
\]

Par suite, \deleted{\unsure{si on}} la relation (*) \deleted{\unsure{équivaut}} \unsure{résultera} de la relation $H^i_x(X(h),\Lambda_\nu)=0$ pour $i\leq m'$, $\forall h$, i.e. \deleted{\unsure{(5.1.3)}} de la relation
\[
\operatorname{prof\,ét}_x(X(h)\geq m'+1\quad\text{pour}\quad\forall\unsure{k|h}.
\tag*{\tikz[baseline=(n.base)]{\node[inner sep=1pt] (n) {5.1.3};\draw (n.center) ellipse (0.62cm and 0.22cm);}}
\]

\unsure{Donc} On \deleted{\illegible{2}} \deleted{\unsure{en}} en conclut aussitôt, par récurrence sur $\dim{\unsure{\mathscr{O}|\mathcal{O}}}_{X_0,x}=r$ :

\uline{\deleted{\unsure{Proposition}}\inserted{Théorème} 5.2.} Soit $Z$ une partie de $X_0$, \inserted{\unsure{n'importe}.} Supposons que : \deleted{pour}

a) pour $\forall$ générisation $x'$ de $x$ dans $X_0$, \deleted{\unsure{distinct de $x$, on}} contenue dans $X_0-Z$, on ait
\[
(\underline{\Phi}^i_\nu)_{\bar{x}'}=0\quad\text{si}\quad i\leq m',
\]
\scanpage{150}
\mspage{V.22}
% archive sheet number: 149

(c'est le cas si on suppose p.ex. que $Z$ est l'ens des \unsure{pts} de $X_0$ en lesquels $f$ n'est pas lisse).

b) Pour toute générisation $x'$ de $x$ dans $Z$, \deleted{\unsure{soit}} \unsure{on|ou} et tt \deleted{\unsure{entier}}\inserted{entier} $\unsure{h|k}$, on a
\[
\operatorname{prof\,ét}_{\bar{x}'}(X(h))\geq m'+1,
\]
\deleted{(ce qui est le cas}
Alors on a
\[
({\unsure{\underline{\Phi}|\underline{\varphi}}}^i_{\unsure{\nu|v}})_x\unsure{=}0\quad\text{si}\quad i\leq m'
\tag*{(5.2.1)}
\]
i.e. $f$ est loc. $(m'\deleted{\unsure{-|+}}1)$-acyclique en $x$, pour \deleted{\unsure{les}} $\{\ell\}$, donc si $f$ est propre, on a
\[
H^i(X_0,\Lambda_\nu)\unsure{-|\longrightarrow}H^i(X_{\bar{\unsure{\eta|y}}},\Lambda_\nu)
\quad\text{un isom si $i<m'$, mon. si $i=m'$.}
\tag*{(5.2.2)}
\]

% The entire block from Corollaire 5.3 to the foot of the page is cancelled.
\deleted{\uline{Corollaire} 5.3. \deleted{Supposons} Soit $Z$ une partie de $X_0$, \deleted{\unsure{soit}}\inserted{$m$} un entier, \deleted{Supposons} posons}
\[
\deleted{\left\{\begin{aligned}
d&=\sup_{\substack{x'\in Z\,\deleted{\unsure{\cap X_{\bar{x}}}}\\x'\text{ générise }x}}\dim\overline{\{x'\}},\\
\deleted{m}&\deleted{=\unsure{\operatorname{prof\,ét}^{*}_Z}}\\
1+m&=\sup_h\unsure{\operatorname{prof\,ét}^{*}_Z(X(h))}\;\unsure{\in|\text{et}}\;{\unsure{\sup|\inf}}_h\unsure{\operatorname{prof\,ét}^{*}(X(h))}-\dim\overline{\{\unsure{x'|x}\}}
\end{aligned}\right.}
\tag*{\deleted{\unsure{(5.3.1)}}}
\]

\deleted{\deleted{Alors on a} Supposons que l'on ait}
\[
\deleted{\deleted{(\underline{\Phi}^i_\nu)_{\bar{x}'}=0\quad\text{si}\quad i\leq}}
\]
\[
\deleted{\text{a)}\quad(\underline{\Phi}^i_\nu)_{\bar{x}'}=0\quad\text{si}\quad i\unsure{\leq|<}m-d\deleted{\unsure{+1}},\quad\text{pour }x'\in X-Z,\quad x\in\overline{\{x'\}},}
\tag*{\deleted{\deleted{(5.3.2)}}}
\]

\deleted{\deleted{Alors on a aussi}}

\deleted{b) Pour toute générisation $x'$ de $x$ dans $X_0$, on a}
\[
\deleted{\deleted{(\underline{\Phi}_{\unsure{\bar{x}|x}})=0\quad\text{si}\quad i\leq m-d}}
\]
\[
\deleted{\operatorname{prof\,ét}^{*}_{x'}(X_0)\geq m\inserted{+1}\quad\text{si }x'\in X_0\quad(\text{resp. }\operatorname{prof\,ét}^{*}_{x'}(X_\eta)=0\quad\text{si }x'\in X_\eta).}
\tag*{\deleted{(5.3.3)}}
\]

\deleted{\deleted{i.e. $f$ est localement $(m-d)$-acyclique en $x$, pour $\{\ell\}$.}}

\deleted{Alors}

\deleted{Résulte \unsure{formellement|facilement} de 5.2, en y faisant $m'=m-d$, la condition a) de 5.2 est satisfaite \unsure{par|en} vertu de (5.3.2), et la condition b) résultant de}
\scanpage{151}
\mspage{V.23}
% archive sheet number: 150

\margin{(Cf. 5.0. sur les hypothèses préliminaires) \deleted{\unsure{Si le cas}}}
\uline{Corollaire} 5.3. \inserted{Supposons la condit. a) de 5.2. satisfaite.} La condition b) \inserted{de 5.2. (et donc la conclusion 5.2.1)} est satisfaite dans chacun des cas suivants (i) et (ii), où on a posé $d=\sup_{\deleted{\unsure{x'\in Z}}}\dim\overline{\{x'\}}$, où $x'$ parcourt les générisations de $x$ dans $X_0$ :

\deleted{\unsure{(i)}} \uline{Cas I} Pour toute générisation $x'$ de $x$ dans $X$, on a
\[
\left\{\begin{aligned}
\operatorname{prof\,ét}^{\unsure{*|\alpha}}_{x'}(X_0)&\geq m'+1&&\text{si }x'\in Z\\
\operatorname{prof\,ét}^{*}_{x'}(X_{\bar{\unsure{\eta|y}}})&\geq m'\deleted{\unsure{+1}}+\dim\overline{\{x\}}&&\text{si }x'\in X_\eta,
\end{aligned}\right.
\tag*{(5.3.1)}
\]
ce qui est le cas si on fait $m'=m-d-1$, avec $m$ un entier \deleted{$m$} tel que
\[
\left\{\begin{aligned}
\operatorname{prof\,ét}^{\unsure{\alpha|*}}_{x'}(X_0)&\geq m&&\text{pour toute \inserted{générisation} $x'\in Z$ de $x$}\\
\operatorname{prof\,ét}^{\alpha}_{x'}(X_{\bar{\eta}})&\geq m\deleted{\unsure{-d}}-1&&\text{pour tt gén. $x'$ de $x$ dans $X$, avec $x'\in X_\eta$.}
\end{aligned}\right.
\tag*{(5.3.2)}
\]

\uline{Cas II} \deleted{Pour toute générisation $x'$ de $x$ dans $Z$,} \inserted{\deleted{On a} Le morphisme $f$ est plat en $x$, et on a} \deleted{on a}
\[
m'=\operatorname{prof\,géom}^{*}_x(X_0)\deleted{\unsure{-d}},\quad\text{où}\quad
\operatorname{prof\,géom}^{*}_x(X_0)\overset{\mathrm{df}}{=}\operatorname{prof\,géom}_x(X_0)+\inserted{\dim\overline{\{x\}}}
\]
Cf. SGA 2 $\overline{\underline{\mathrm{XIV}}}$ 5.3 pour la notion de profondeur géométrique). \deleted{ce qui est le cas si on fait $m'=m-d$, avec $m$ un entier tel que}
\[
\deleted{\operatorname{prof\,géom}^{*}_x(X)}
\]

\uline{Démonstration}

\uline{Cas I} Les relations (5.3.1) étant stables par \inserted{générisation et} changement de $X$ en $X(\unsure{h|k})$, il suffit de vérifier qu'elles \deleted{\illegible{2}} entraînent
\[
\operatorname{prof\,ét}_x(X)\geq m'+1.
\tag*{(*)}
\]
Or la deuxième des relations (5.3.1) \deleted{\unsure{peut aussi}} \inserted{s'écrit} \deleted{s'écrire}
\[
\operatorname{prof\,ét}_{x'}(X')\deleted{\unsure{\geq m'}}+\dim\overline{\{x'\}}^{(\unsure{X_\eta|X})}\geq m'+\dim\overline{\{x\}},
\]
\margin{$=\dim\overline{\{x\}}+\dim\overline{\{\bar{x}'\}}^{(X')}$}
or
\[
\dim\overline{\{x'\}}^{(X_\eta)}=\dim\overline{\{x'\}}^{X}-1=\dim\overline{\{x\}}+\dim\overline{\{\bar{x}'\}}^{(X')}-1
\]
où $\bar{x}'$ est un pt de $X'$ au dessus de $x'$, donc on \unsure{peut|voit}
\scanpage{152}
\mspage{V.24}
% archive sheet number: 151

l'\unsure{interpréter} en disant que pour $x'\in X'_{\unsure{\eta|y}}=X'-X'_0$, on a
\[
\operatorname{prof\,ét}_{x'}(\deleted{\unsure{X'}}V)+\dim\overline{\{x'\}}^{(V)}\geq m'\deleted{\unsure{+1+\dim\overline{\{x\}}}},
\]
ce qui par le th. de Lefschetz local SGA 1 $\overline{\underline{\mathrm{XIV}}}$ 4.5 implique que
\[
H^i(\deleted{\unsure{X'}}V)\longrightarrow H^i(V_s)\quad\text{est bijectif pour $i<m'-1$,}
\]
\inserted{injectif pour $i=m'-1$}
% slanted
\margin{il \unsure{suffirait} \unsure{que ce soit} bijectif pour $i<m'-1$, injectif pour $i=m'-1$, sauf dans le cas où $m'=1$}

donc \deleted{il suffit pour établir} \inserted{la relation} (*) \deleted{de prouver que l'on a} équivaut à $\operatorname{prof\,ét}_x(X_0)\geq m'+1$, ce qui est la première des \deleted{\unsure{prof}} relations (5.3.1). Le fait que les relations (\unsure{5.3.2|5.3.8}) impliquent les relations (5.3.1) est d'autre part immédiat.

\uline{Cas II} De même, l'hypothèse est \inserted{évidemment} stable par passage de $X$ à $X(\unsure{h|k})$, \deleted{de plus elle est stable par générisation, comme on constate aussitôt que} \inserted{Il suffit donc de prouver} qu'elle implique (*)
\[
\deleted{\operatorname{prof\,ét}_{x'}(X)\geq m'+1}
\tag*{\unsure{(**)}}
\]
pour tt gén. $x'$ de $x$ dans $Z$. Or on a (lemme \deleted{à dégager !}\inserted{5.3.3 (i) (plus bas)})
\[
\deleted{\operatorname{prof\,ét}_{x'}(X)\geq\operatorname{prof\,ét}_{x'}}
\]
\[
\operatorname{prof\,géom}_{x'}(X_0)\geq\operatorname{prof\,géom}_{x}(X)\unsure{-|\deleted{+}}\operatorname{codim}(\bar{x},\bar{x}')
\]
donc, comme $\operatorname{codim}(\bar{x},\bar{x}')=\dim\overline{\{x'\}}-\dim\overline{\{x\}}\unsure{\geq|\leq}d-\dim\overline{\{x\}}$ on trouve
\[
\operatorname{prof\,géom}_{x'}(X_0)\geq\operatorname{prof\,géom}^{*}_x(X_0)\deleted{\unsure{+\dim\overline{\{x\}}}}-d=m'.
\]
D'autre part, (lemme 5.3.3 (ii) (plus bas)) on a
\[
\operatorname{prof\,géom}_{x'}(X)\geq\operatorname{prof\,géom}_{x'}(X_0)+\underset{\substack{\Vert\\\operatorname{prof\,géom}_s(S)}}{1},
\]
donc on en conclut
\[
\operatorname{prof\,géom}_{x'}(X)\geq m'+1.
\]
\scanpage{153}
\mspage{V.25}
% archive sheet number: 152

D'autre part, le th de Lefschetz local entraîne (SGA 1 $\overline{\underline{\mathrm{XIV}}}$ 5.6)
\[
\operatorname{prof\,ét}_{x'}(X)\geq\operatorname{prof\,géom}_{x'}(X),
\]
d'où \deleted{par} (**) , cqfd.

\unsure{Il reste à regarder un}

\uline{Lemme} 5.3.3 (i) Soit $A$ un anneau local noethérien, $\unsure{\mathfrak{p}|p}$ un idéal premier de $A$. Alors on a
\[
\operatorname{prof\,géom}A_{\mathfrak{p}}\geq\operatorname{prof\,géom}A-\dim A/\mathfrak{p}
\]

(ii) Soit $A\unsure{-|\longrightarrow}B$ un hom local d'anneaux locaux noethériens, avec $B$ \uline{plat} sur $A$, $k$ le corps résiduel de $A$, $K$ celui de $B$, et supposons qu'il existe un anneau local noethérien $C$ \deleted{sur $k$, de corps} formellement lisse sur $k$ de corps résiduel $K$ (*) (condition satisfaite en tous cas si $K$ est une extension de type fini de $k$). Alors on a
\[
\operatorname{prof\,géom}B\geq\operatorname{prof\,géom}A+\operatorname{prof\,géom}B\otimes_A k.
\]

Donner \unsure{démonstrations|démonstration}, !

(*) Cette condition est peut-être inutile. Pour une reformulation de cette condition, cf EGA $0_{\mathrm{IV}}$ 22.2.2
\scanpage{154}
\mspage{V.26}
% archive sheet number: 153

Voici un corollaire du Cas II :

\uline{Corollaire} 5.4 Supposons \inserted{$X$ plat sur $S$ et} $X_0$ loc. une intersection complète et purement de dimension $n$. Alors on a
\[
{\unsure{\underline{\Phi}|\underline{\varphi}}}^i_{\unsure{\nu|v}}=0\quad\text{pour}\quad
\deleted{i\notin[n-d+1,n+1]\quad\unsure{i\leq m'}}\inserted{i\notin[n-d+1,n+1]},
\]
où $d$ est la dimension de l'ens $Z$ des pts de $X_0$ en lesquels $f$ est non lisse, \deleted{donc} et
\[
\Phi^i_\nu=0\quad\text{si}\quad i\notin[n-d+1,n+d+1]
\]
En particulier, si $d=0$, on a
\[
\begin{aligned}
\underline{\Phi}^i_\nu&=0\quad\text{si}\quad i\neq n+1\\
\Phi_\nu&=0\quad\text{si}\quad i\neq n+1,
\end{aligned}
\]
donc l'homomorphisme
\[
H^i(X_0,\Lambda_\nu)\longrightarrow H^i(X_{\bar{\unsure{\eta|y}}},\Lambda_\nu,
\]
est un isomorphisme si $i\neq n,n+1$, un monomorphisme pour $i=n$, un épimorphisme pour $i=n+1$; \inserted{de plus} L'action du grp d'inertie \deleted{\unsure{de}} $I$ sur $H^i(X_{\bar{\eta}},\Lambda_\nu)$ \deleted{est triviale} ou sur $H^i(X_{\bar{\eta}},\mathbb{Z}_\ell)$ est triviale si $i\neq n$.

5.5. Pour terminer le présent paragraphe, nous allons donner une \unsure{variante} des résultats précédents pour la cohomologie non commutative. Supposons $f$ loc. $1$-acyclique \inserted{pour $\unsure{\mathbb{L}|L}$} en toutes les générisations de $x$, \inserted{distinctes de $x$} et \unsure{proposons-nous|proposons} de \deleted{donner} trouver des conditions de \unsure{locale|local} $1$-asphéricité en $x$, pour l'ens de nbs premiers \deleted{$\unsure{\ell\in\mathbb{P}-\{p\}}$}.
\diagram{\begin{tikzcd}[ampersand replacement=\&]
V \arrow[r,no head] \arrow[d,phantom,"\cap" description] \& \bar{V}\arrow[dd,no head]\\
X'\arrow[d] \& \\
S \& \unsure{\bar{S}|\bar{\eta}}\arrow[l]
\end{tikzcd}}
Je dis que pour tt \inserted{$\mathbb{L}$-}grp fini $G$, \deleted{premier à $\mathbb{L}$}, on a
\[
H^i(\bar{V},G)\xrightarrow{\sim}H^i(V_{\bar{\eta}},G)\quad\text{si}\quad i\leq1.
\tag*{(5.5.1)}
\]
\scanpage{155}
\mspage{V.27}
% archive sheet number: 154

Donc pour avoir la locale $1$-asphéricité pour $\unsure{\mathbb{L}|L}$ en $x$, il suffit de supposer que l'on a, pour tt tel $G$ :
\[
H^i(\deleted{\unsure{X}}V(\unsure{h|k}),G)=0\quad\text{si}\quad i\unsure{\leq|<}1
\tag*{\unsure{(5.5.1)|(5.5.2)}}
\]
On sait \deleted{\unsure{par suite de|par vertu de}} que cette condition est vérifiée \unsure{les|lors} \deleted{\unsure{pour|que}} $\unsure{\deleted{X(h)}|X(h)}$ \deleted{\unsure{est|sont}} \unsure{des|les} intersections complètes de dim $\geq3$ en $x$, \inserted{(et sans doute dès que leur prof géom en $x$ est $\geq3$)}, lorsque $f$ est plat en $x$, cela signifie aussi que $X_0$ est une int. compl. de dimension $\geq2$ en $x$.

Utilisant le th de locale $1$-asphéricité pour $\mathbb{L}$ des morphismes lisses, on trouve ainsi le

\uline{Théorème} 5.6. Supposons que $f:X\longrightarrow S$ soit plat, que l'ens $Z$ des points \inserted{de $X$} en lesquels $f$ est non lisse est de codimension $\geq2$, et que $X_0$ soit une intersection complète locale en tous ses points. Alors $f$ est localement $1$-asphérique pour $\mathbb{L}_{\unsure{\eta|s}}=\mathbb{P}-\{p\}$. En particulier si $f$ est propre, \inserted{alors,} pour $\forall$ groupe fini $G$ \inserted{d'ordre} premier à $p$, l'hom de spécialisation
\[
H^{\unsure{1|i}}(X_0,G)\unsure{-|\longrightarrow}H^1(X_{\bar{\unsure{\eta|y}}},G)
\]
est un isomorphisme.

\uline{Remarque} 5.7. a) Il semble qu'il doive être possible de remplacer l'hypothèse que $X_0$ est une intersection complète par celle que pour tt $x\in Z$, on ait $\operatorname{prof\,géom}_x(X_0)\geq2$; cela résulterait d'une \unsure{solution|réponse} affirmative à SGA 1 $\overline{\underline{\mathrm{XIV}}}$ 6.3
\scanpage{156}
\mspage{V.28}
% archive sheet number: 155

b) Les résultats précédents suggèrent la conjecture suivante \inserted{(au moins pour $A,B$ \uline{excellents})}

Soit $f:X\inserted{=\operatorname{Spec}B}\unsure{-|\longrightarrow}S\inserted{=\operatorname{Spec}A}$ un morphisme \uline{plat} de schémas noethériens, \inserted{stricts locaux, $s$ le point fermé de $S$. $X_0=f^{-1}(s)$.} \deleted{\unsure{et à les|et à la} fibres géométriques \unsure{satisfaisant|satisfassant} à la condition $(R_d)$,} \deleted{Soit \deleted{$X_0$} $k$ le corps \unsure{rés. du pt} \illegible{1},} Supposons que \deleted{$X_0$ soit géométriquement $(R_d)$, et que} pour tt extension finie $k'$ de $k=k(s)$, l'ens. \deleted{$m=\operatorname{prof\,géom}$}\inserted{singulier de} $X_0\otimes_k k'$ \deleted{satisfasse la condition $R_{\unsure{1|d}}$, et supposons} \inserted{soit de dim $\leq d$} et $m=\operatorname{prof\,géom}B$. Alors $f$ est loc. $(m-d\unsure{-|+}1)$-acyclique, \deleted{et si $m-d\unsure{-|+}1\geq\deleted{\unsure{3}}1$} pour $\unsure{\mathbb{L}|L}=\mathbb{P}-\{p\}$, ($p$. car résiduelle de $A$), et si $m-d-1\geq1$ i.e. $m-d\geq2$, $f$ est loc. $1$-asphérique pour $\mathbb{L}$. \deleted{Cela} Il est plausible que cela soit démontrable comme conséquence de la conjecture déjà citée SGA 1 $\overline{\underline{\mathrm{XIV}}}$ 6.3, et de la résolution des singularités pour les schémas \unsure{finis} sur $\unsure{S|\eta}$,
\scanpage{157}
\mspage{V.29}
% archive sheet number: 156

6. \uline{Cas d'un point de non lissité \inserted{dans $X_0$} maximal}. \deleted{\unsure{de $X_0$}}

(6.1) Soit $x\in X_0$, et \deleted{suppo} reprenons les notations du § 4. Nous supposons que l'on a pour tt $i$
\[
\tag*{(6.1.1)}
{\unsure{\underline{\Phi}|\underline{\varphi}}}^i_{\unsure{\nu|v}}|V=0,\quad\text{i.e. }(\underline{\Phi}^i_\nu)_{\deleted{\unsure{\bar{x}}}\bar{x}'}=0\quad\text{pour tte gén. }x'\ne x\text{ dans }X_0\ \deleted{\unsure{en}}.
\]
Condition satisfaite en particulier si $x$ est un point maximal de l'ensemble $Z$ des pts de $X_0$ en lesquels $f$ n'est pas lisse, \deleted{\unsure{et lorsque plus il en est|et lorsque plus il a} (4.4.3) \unsure{une comp.}\ \unsure{conn.}\ point de $X$ \unsure{dominant et contenant} $x$}.

L'inclusion
\[
V_{\bar{\unsure{\eta|y}}}\longleftarrow\bar V
\]
définit donc un isomorphisme
\[
\tag*{(6.1.2)}
\mathbb{R}\unsure{\Gamma|\mathscr{T}}(\bar V,\Lambda_\nu)\xrightarrow{\sim}\mathbb{R}\Gamma(V_{\bar\eta},\Lambda_\nu)
\]
i.e. des isomorphismes
\[
\tag*{(6.1.3)}H^i(\bar V,\Lambda_\nu)\simeq H^i(V_{\bar\eta},\Lambda_\nu),
\]
On peut \unsure{en sus} \unsure{compter|considérer} à opérateurs $\unsure{\pi|\Pi}$, (6.1.2) \deleted{\unsure{donne}} s'écrit
\[
\tag*{(6.1.4)}\mathbb{R}(\Gamma,\pi)(\bar V,\Lambda_\nu)\xrightarrow{\sim}\mathbb{R}(\Gamma,\pi)(V_{\bar\eta},\Lambda_\nu\unsure{,|)}.
\]
Donc, au point de vue de leur structure \unsure{chomologique|cohomologique}, la situation des 5 schémas (4.2) se réduit aux quatre schémas $V_s\simeq V_{\bar s},V,V_\eta,V_{\bar\eta}$.

(6.2) L'inclusion $V_{\bar s}\longrightarrow\bar V$ donne donc un homomorphisme \inserted{${\unsure{K|\kappa}}^{\unsure{\bullet|\prime}}$}
\inserted{($(L^\bullet)^{\mathrm{triv}}$, où $L^\bullet=\mathbb{R}\Gamma(V_s,\Lambda_\nu)$)}
\[
\tag*{(6.2.1)}
\boxed{u=(\kappa_\pi^{\unsure{*|\bullet}})^{\unsure{*|\bullet}}:\mathbb{R}(\Gamma,\pi)(V_{\bar\eta},\Lambda_\nu)\longrightarrow\mathbb{R}(\Gamma,\pi)(V_{\bar s},\Lambda_\nu),}
\]
\margin{compte tenu que $\pi$ opère trivialement sur $V_{\bar s}=V_s$}
qui \deleted{par le} s'insère dans le triangle distingué
\diagram{\[
\tag*{(6.2.2)}
\begin{tikzcd}[ampersand replacement=\&,column sep=small]
\& {\mathbb{R}(\Gamma,\pi)(V_{\bar s},\Lambda_\nu)\simeq\mathbb{R}\Gamma(V_s,\Lambda_\nu)^{\mathrm{triv}}=L^{\bullet\,\mathrm{triv}}}\arrow[dl] \& \\
{R^\bullet=\mathbb{R}(\Gamma,\pi)(\bar V,(\Lambda_\nu)_{V_{\bar\eta},\bar V})}\arrow[rr] \& \& {\mathbb{R}(\Gamma,\pi)(V_{\bar\eta},\Lambda_\nu)=K^\bullet}\arrow[ul]
\end{tikzcd}
\]}
\scanpage{158}
\mspage{V.30}
% archive sheet number: 157

L'homomorphisme $u$, (\deleted{et} \inserted{ou ce qui revient au même,} le triangle (6.2.2) correspondant) domine l'étude cohomologique des quatre schémas $V_s\simeq V_{\bar s},V,V_{\unsure{\eta|y}},V_{\bar\eta}$, en ce sens qu'il permet de retrouver formellement la cohomologie (à coefficients constants) des schémas envisagés. Appliquant le foncteur ``oubli de $\unsure{\pi|\Pi}$'' à 6.2.2, on trouve le triangle exact
\diagram{\[
\tag*{(6.2.3)}
\begin{tikzcd}[ampersand replacement=\&,column sep=small]
\& {\mathbb{R}\unsure{\Gamma|\mathscr{T}}(V_s,\Lambda_{\unsure{\nu|v}})=L^\bullet}\arrow[dl] \& \\
{R^{\bullet\,\mathrm{oub}}=\mathbb{R}\Gamma(\bar V,\Lambda_{\nu\unsure{,|\,}V_{\bar\eta},\bar V})}\arrow[rr] \& \& {\mathbb{R}\Gamma(V_{\bar\eta},\Lambda_\nu)=K^{\bullet\,\mathrm{oub}}}\arrow[ul,"{\unsure{\kappa|k|h}}^*"]
\end{tikzcd}
\]}
qui donne la cohomologie (sans opérateurs) de $V_{\bar\eta}$ et $V_s$ et l'hom du premier dans le second défini par l'inclusion
\diagram{\[
\begin{tikzcd}[ampersand replacement=\&]
V_{\bar s}\arrow[r,"\bar i"]\arrow[rr,bend right=25,dashed,"\kappa"'] \& \bar V \& V_{\bar\eta}\arrow[l,"\sim"',"\mathrm{coh.}"]
\end{tikzcd}
\]}.
D'autre part, \unsure{appliquons|appliquant} le foncteur $\mathbb{R}\Gamma^\pi$ \inserted{à (6.2.1), (6.2.2)}, on trouve
\diagram{\[
\tag*{(6.2.4)}
\begin{tikzcd}[ampersand replacement=\&,column sep=small]
\& {\mathbb{R}\Gamma(V_s,\Lambda_\nu)\mathbin{\mathop{\otimes}^{\mathrm L}}_{\Lambda_\nu}\mathbb{R}\Gamma^\pi(\Lambda_\nu)\simeq\deleted{\unsure{\mathbb{R}\Gamma}}L^\bullet\oplus L^\bullet[-1]}\arrow[dl] \& \\
{\mathbb{R}\Gamma^\pi(R^\bullet)}\arrow[rr] \& \& {\mathbb{R}\Gamma(V_\eta,\Lambda_\nu)\simeq\mathbb{R}\Gamma^\pi(K^\bullet)}\arrow[ul,"\mathbb{R}\Gamma^\pi(\kappa^*)\unsure{=|\sim}"]
\end{tikzcd}
\]}
\inserted{\unsure{dégager la formule|dégén. formule}}
% margin note written slanted
\margin{\unsure{remonter|remettre} \unsure{ce|le} diagramme et le commenter !
\diagram{\begin{tikzcd}[ampersand replacement=\&]
V \& V_\eta\arrow[l,"\deleted{\unsure{g}}\unsure{j_s|j_\eta|j}"]\\
{V_s\simeq V_{\bar s}}\arrow[u,"\deleted{\unsure{h_s}}i"]\arrow[r,dashed,"\kappa"] \& {\bar V\overset{\deleted{\unsure{\mathrm{coh.}}}}\simeq V_{\bar\eta}}\arrow[u,"h_\eta\deleted{\unsure{g}}"']
\end{tikzcd}}
(avec des notations légèrement changées p.r. à (4.2))}
compte tenu que, par la structure \unsure{spéciale|précise} de $\pi=I$, on a
\[
\tag*{(6.2.5)}\mathbb{R}\Gamma^\pi(\Lambda_\nu)\simeq\Lambda_\nu\oplus\underbrace{\Lambda_\nu(-1)}_{\text{twist de Tate}}\overbrace{[-1]}^{\text{transl. des degrés}}.
\]
\deleted{L'homomorphisme $\mathbb{R}\Gamma^\pi(\kappa^*)$ \unsure{défini|définit} se décompose en
$\mathbb{R}\Gamma(V_\eta,\Lambda_\nu)\longrightarrow$}
Ainsi on retrouve $\mathbb{R}\Gamma(V_\eta,\Lambda_\nu)$, et l'hom $h^*$ déduit de $V_{\bar\eta}\xrightarrow{h}V_\eta$ n'est autre que l'hom naturel ``d'inclusion'' $\mathbb{R}\Gamma^\pi(K^\bullet)\longrightarrow\deleted{\unsure{\mathbb{R}\Gamma}}\{K^\bullet\}$. De plus, l'hom $\mathbb{R}\Gamma^\pi(\kappa^*)$ se décompose en
\[
\tag*{(6.2.6)}\mathbb{R}\Gamma^\pi(\kappa^*)=(\alpha,\beta):\mathbb{R}\Gamma(V_\eta,\Lambda_\nu)=\mathbb{R}\Gamma^\pi(K^\bullet)\longrightarrow\mathbb{R}\Gamma(V_s,\Lambda_\nu)+\mathbb{R}\Gamma(V_s,\Lambda_\nu)\unsure{[-1]|(-1)}
\]
\scanpage{159}
\mspage{V.31}
% archive sheet number: 158

où $\alpha$ \deleted{\unsure{est un}} est donné par
\[
\tag*{(6.2.7)}\alpha={\unsure{\kappa|k}}^*h^*,
\]
tandis que $\beta$ est un hom. de degré 1 :
\[
\tag*{(6.2.8)}\unsure{\beta|\rho}:\mathbb{R}\unsure{\Gamma|\mathscr{T}}(V_{\unsure{\eta|y}},\Lambda_{\unsure{\nu|v}})\longrightarrow\mathbb{R}\Gamma(V_s,\Lambda_\nu)(-1)[-1].
\]
\deleted{\uline{\unsure{Proposition}}} Cet homomorphisme peut s'interpréter, par la proposition suivante, qui en \unsure{même temps|toujours} donne un \deleted{\unsure{avec|autre} \unsure{un mode de}} calcul de $\mathbb{R}\Gamma(V,\Lambda_\nu)$, en termes de la donnée fondamentale \deleted{\unsure{en u|en de}} (6.2.1) :

\uline{Proposition} 6.3 a) On a
\deleted{$\underline{H}^i_{V_s}(\unsure{V_s|V},\Lambda_{\nu V})=0$ pour $i\ne2$, et $\underline{H}^2_{V_s}(\Lambda_{\nu V})\simeq\Lambda_{\nu V_s}(1)$,}
\[
\tag*{(6.3.1)}\underline{H}^i_{V_s}(\Lambda_{\nu V})\simeq\begin{cases}
0&\text{si }i\ne2\\
\deleted{\text{si }i}&\\
\Lambda_{\nu V_s}(-1)&\text{si }i=2,
\end{cases}
\]
en d'autres termes :
\[
\tag*{(6.3.2)}\mathbb{R}\underline\Gamma_{V_s}(\Lambda_{\nu V})\simeq\Lambda_{\nu V_s}(-1)[-2],
\]
d'où un isom
\[
\tag*{(6.3.3)}\mathbb{R}\Gamma_{V_s}(V,\Lambda_\nu)\simeq\mathbb{R}\Gamma(V_s,\Lambda_\nu)(-1)[-2]\inserted{\text{ can.}}
\]
b) Considérons le triangle canonique
\diagram{\[
\tag*{(6.3.4)}
\begin{tikzcd}[ampersand replacement=\&,column sep=small]
\& {\mathbb{R}\Gamma(V_\eta,\Lambda_\nu)}\arrow[dl,"\partial"'] \& \\
{\mathbb{R}\Gamma(V_s,\Lambda_\nu)\unsure{(-1)[-2]|[1][-2]}\underset{\text{par (\unsure{6.3.8|6.3.3})}}\simeq\mathbb{R}\Gamma_{V_s}(\deleted{\unsure{V_s}}V,\Lambda_{\nu V})}\arrow[rr,"i_*"] \& \& {\mathbb{R}\Gamma(V,\Lambda_\nu)}\arrow[ul,"j^*"]
\end{tikzcd}
\]}.
Alors la flèche $\partial$ du dit triangle s'identifie [via l'isomorphisme (6.3.3)] à $\beta$. \deleted{Donc $\mathbb{R}\Gamma(V,\Lambda_\nu)$ \unsure{s'identifie à} \unsure{isomorphe}}
\margin{signe ?}
\deleted{$\mathbb{R}\Gamma(V,\Lambda_\nu)$}

\uline{Corollaire} 6.4 \inserted{À des signes et shifts près,} Le triangle (6.3.4) n'est autre que le triangle distingué \deleted{construit avec} \inserted{des} ``{\unsure{mapping-cylindres}}'' du $\beta$ de (6.2.8).
\scanpage{160}
\mspage{V.32}
% archive sheet number: 159

\uline{Démonstration} le a.

% The following diagram and draft calculation are cancelled together.
\diagram{\begin{tikzcd}[ampersand replacement=\&]
\deleted{V_{\bar{\unsure{\eta|y}}}}\arrow[d,"\deleted{h}"']\arrow[dr,"\deleted{g}"] \& \\
\deleted{V_\eta}\arrow[r,"\deleted{j}"'] \& \deleted{V}
\end{tikzcd}}
\[
\deleted{\mathbb{R}g_*(\Lambda_{V_{\bar\eta}})\unsure{\simeq|\ne}\mathbb{R}j_*\overbrace{\mathbb{R}h_*(\Lambda_{V_{\bar\eta}})}^{h_*(\Lambda_{V_{\bar\eta}})}}
\]
\deleted{et induisant sur $V_s\unsure{\leftarrow|\hookleftarrow}V$, et prenant les invariants on a}
\[
\deleted{\mathbb{R}g_*(\Lambda_{V_{\bar\eta}})\simeq\mathbb{R}j_*\mathbb{R}h_*(}
\]

\diagram{\begin{tikzcd}[ampersand replacement=\&]
V_{\bar s}\arrow[r,"\bar i"]\arrow[d,"g_0"'] \& \bar V\arrow[d,"g"] \& V_{\bar\eta}\arrow[l,"\bar j"']\arrow[d,"h"]\\
V_s\arrow[r,"i"'] \& V \& V_\eta\arrow[l,"j"]
\end{tikzcd}}
\[
\begin{aligned}
\mathbb{R}j_*(\Lambda_{V_\eta})&\simeq\mathbb{R}j_*\mathbb{R}{\unsure{\Gamma|\mathscr{T}}}^{\unsure{\pi|\Pi}}\mathbb{R}h_*(\Lambda_{V_{\bar\eta}})\\
&\simeq\mathbb{R}\Gamma^\pi\mathbb{R}j_*\mathbb{R}h_*(\Lambda_{V_{\bar\eta}})\\
&\simeq\mathbb{R}\Gamma^\pi\mathbb{R}g_*\mathbb{R}\bar j_*(\Lambda_{V_{\bar\eta}})\\
&\xrightarrow{\sim}\mathbb{R}\Gamma^\pi\mathbb{R}g_*(\Lambda_{\bar V})\qquad\text{car }\Lambda_{\bar V}\simeq\mathbb{R}\bar j_*\Lambda_{V_{\bar\eta}}
\end{aligned}
\]
d'où
\[
\begin{aligned}
i^*\mathbb{R}j_*(\unsure{V_{V_\eta}|\Lambda_{V_\eta}})&\simeq\mathbb{R}\underline\Gamma^\pi i^*\mathbb{R}g_*(\Lambda_{\bar V})\\
&\simeq\mathbb{R}\underline\Gamma^\pi\mathbb{R}g_{0*}(\Lambda_{V_{\bar s}})\\
&\simeq\mathbb{R}\underline\Gamma^\pi(\Lambda_{V_s})^{\mathrm{triv}}\\
&\simeq\Lambda_{V_s}\unsure{-|\,}\oplus\Lambda_{V_s}(-1)[-1]
\end{aligned}
\]
d'où aussitôt 6.3. a), en utilisant les relations habituelles entre les $\underline H^i_{V_s}(\Lambda_V)$ et les $\mathbb{R}^ij_*(\Lambda_{V_\eta})$.

Compatibilité b) laissée \unsure{aux rédacteurs|au rédacteur} en forme !
\scanpage{161}
\mspage{V.33}
% archive sheet number: 160

6.5. Utilisons maintenant les résultats du § 2, qui \unsure{impliquent|indiquent}, \deleted{\unsure{si} $n=\dim{\unsure{\mathscr{O}|\mathcal{O}}}_{X_0,x}$} pour
\[
\tag*{(6.5.1)}n=\dim\mathscr{O}_{X_0,x}
\]
que l'on a
\[
\tag*{(6.5.2)}H^i(V_{\bar{\unsure{\eta|y}}})\deleted{\unsure{\Lambda|\Lambda_\nu}}=0\quad\text{si }i>n.
\]
\deleted{\unsure{Complétons}}\inserted{Utilisons les} \unsure{deux|des} \deleted{la} suites exactes déduites de \deleted{6.3.4} la suite spectrale de Hochschild-Serre
\[
\deleted{H^i(V)\unsure{-|\longrightarrow}H^i(V_{\bar\eta}}
\]
\[
H^{\unsure{*|t}}(V_\eta)\Longleftarrow E_2^{pq}:H^p(\unsure{\pi|\Pi},H^q(V_{\bar\eta}))
\]
[déduite de $\mathbb{R}\unsure{\Gamma|\mathscr{T}}(V_\eta)\simeq\mathbb{R}\Gamma^\pi\mathbb{R}\Gamma(V_{\bar\eta})$] :
\[
\tag*{(6.5.3)}\unsure{0|:}\unsure{-|\longrightarrow}H^{i\unsure{-|+}1}(V_{\bar\eta})_\pi\longrightarrow H^i(V_\eta)\unsure{-|\longrightarrow}H^i(V_{\bar\eta})^\pi\longrightarrow0,
\]
on trouve donc
\[
\tag*{(6.5.4)}H^i(V_\eta)=0\quad\text{si }i>n+1.
\]
\deleted{Utilisons} Écrivons la suite exacte déduite de 6.3.4
\[
\tag*{(6.5.5)}\longrightarrow H^{i-\unsure{2|1}}(V_s)(-1)\deleted{\unsure{\longrightarrow H^i(V)}}\longrightarrow H^i(V)\longrightarrow H^i(V_\eta)\longrightarrow H^{i-1}(V_s)\longrightarrow H^{i\unsure{+|-}1}(V)\longrightarrow\cdots,
\]
on \deleted{trouve} voit que 6.5.4 s'exprime aussi par \unsure{la relation|les relations} :
\[
\tag*{(6.5.6)}\begin{aligned}
i_*:H^{\deleted{\unsure{i-2}}i}\unsure{(V_s,\Lambda)|(V_s)(-1)}&\longrightarrow H^{i+2}(V_{\deleted{\unsure{s|\eta}}})\quad\text{un isom. si }i\deleted{\unsure{\geq}}>n\deleted{\unsure{-1|+1}}\\
&\qquad\text{un épimorph. si }i=n\deleted{\unsure{-1|+1}}
\end{aligned}
\]

\unsure{6}
\scanpage{162}
\mspage{V.34}
% archive sheet number: 161

6.6. Utilisant les résultats du § 5, on trouve
\[
\tag*{(6.6.1)}\left\{\begin{aligned}
&H^i(\unsure{V.|V})\longrightarrow\underset{\inserted{\begin{gathered}\unsure{\wr|\Vert}\\(\varinjlim H^i(V(m)))\end{gathered}}}{H^i(V_{\bar{\unsure{\eta|y}}})}\longrightarrow H^i(V_s)\qquad\begin{array}{l}\text{des isom si }i\leq n-2\\\text{des mon si }i=n-1\end{array}\\
&\text{\textcircled{si}}\ X\text{ est une intersection complète en toute générisation}\ \deleted{\unsure{stricte}~\unsure{x}\ \unsure{X_0}}\\
&\qquad\inserted{x'\text{ de }x\text{ dans }V-V_s=V_\eta}\ \deleted{\unsure{X_0}}\\
&\text{et si }X\text{ est équidimensionnel (de dim }n+1)\text{ en }x
\end{aligned}\right.
\]
\[
\tag*{(6.6.2)}\left\{\begin{aligned}
H^i(V_s)&=0\quad\text{si }0<i\leq n-2\qquad(H^0(V_s)\simeq\Lambda\text{ si }n\geq2)\\
H^i(V)&=0\quad\text{si }0<i\leq n-1\qquad(H^0(V_s)\simeq\Lambda\text{ si }n\geq1)\\
H^i(V_{\bar\eta})&=0\quad\text{si }i\ne0,n\\
H^i(V_\eta)&=0\quad\text{si }i\ne0,1,n,n+1
\end{aligned}\right|\quad\text{\textcircled{si}}\ \unsure{X_0|X}\text{ intersection complète en }x
\]

\deleted{Supposons}
6.7. Supposons \deleted{enfin} $x$ point maximal de l'ens des points de non lissité de $f$ \uline{dans $X$}, donc à fortiori $X$ régulier en toute générisation stricte de $x$. Alors $H^{\unsure{q|\bullet}}(V_s)$ et $H^q(V)$ satisfont \deleted{à la dualité} \inserted{\unsure{sont des espaces|ont des espaces}} de Poincaré, \deleted{\unsure{de dim}}\inserted{\unsure{pour les dim}} \deleted{en dim} $2n-1$ et $2n+1$ resp ; Par suite si $X$ est \inserted{de plus} une int. compl. en $x$, on déduit de 6.6.2 ; \deleted{\unsure{on trouve}}
\[
\tag*{(6.7.1)}\left\{\begin{aligned}
H^i(V_s)&=0\quad\text{si }i\ne0,n-1,n,2n-1\\
H^i(V)&=0\quad\text{si }i\ne0,n,n+1,2n+1
\end{aligned}\right.
\]
\inserted{$\left[\begin{aligned}&\text{et }H^0(V_s)\simeq\Lambda\ \deleted{\text{si }\unsure{n\geq2}}\\&H^{2n-1}(V_s)\simeq\Lambda(-n\unsure{+|-}1)\ \deleted{\unsure{\text{et}}}\end{aligned}\middle|\text{si }n\geq2\right.$}
\inserted{$\left[\begin{aligned}&\text{et }H^0(V)\simeq\Lambda\ \deleted{\text{si }\unsure{n\geq1}}\\&H^{2n+1}(V)\simeq\Lambda(-n-1)\end{aligned}\middle|\text{si }n\geq1\right]$}
\scanpage{164}
\mspage{1}
% archive sheet number: 163

\unsure{VI|VII} \uline{Groupe de monodromie et groupe fondamental} $(*)$.

\deleted{\illegible{1}} Dans les exposés précédents, nous avons étudié, \deleted{l'ac} pour un schéma de type fini $X$ sur $K$ (corps des fractions de l'anneau de val. \unsure{discr.|discr} str. local $V$), l'action du groupe d'inertie $I=\operatorname{Gal}(\bar K/K)$ sur les gpes de cohomologie de $X_{\bar K}$. Il y a lieu également d'étudier l'action de $I$ sur \deleted{les gpes} d'autres invariants de nature topologique \deleted{\unsure{associés}} de $X_{\bar K}$. Dans le présent exposé, nous nous intéresserons surtout, à l'action de $I$ sur le groupe fondamental ${\unsure{\pi|\Pi}}_1(X_{\bar K},\unsure{\xi|\zeta})$ (relativement à un pt géom $\xi$ de $\bar K$). \deleted{Nous désignerons par $\pi_1^{(\unsure{p|p'})}(X_K,\xi)$ le}

1. \deleted{\unsure{Résultats}}\uline{Énoncé des résultats}.

1.1. Pour un groupe profini $\pi$, nous désignerons par la suite par $\pi^{(p')}$ le plus grand quotient pro-fini de $\pi$ premier à $p$ \inserted{(\unsure{i.e.|soit} $p'=\mathbb{P}-\{p\}$, et)} (où $p$ est l'exp. car. de $k=k(V)$). Donc si $p=1$, \inserted{i.e. $k$ de car. nulle,} on a $\pi^{(p')}=\pi$, sinon $\pi^{(p')}$ est le quotient de $\pi$ par le ss-gpe invariant fermé \uline{engendré} par \unsure{un|son} $p$-ss-groupe de Sylow, i.e. le quotient de $\pi$ par le ss-gpe engendré par les $p$-ss-groupes de Sylow de $\pi$.
% margin note written sideways
\margin{$(*)$ Le présent exposé n'est pas nécessaire pour la compréhension du reste du séminaire.}

La formation de $\pi^{(p')}$ est évidemment fonctorielle en $\pi$, et transforme autos intérieurs de $\pi$
